Das Volumen des Zylinders ist $V = \pi r^2 h$, seine Dichte also
\begin{align}
 \rho &= \frac{m}{V} = \rhoZF \\
   &= \frac{\mZy}{\pi\cdot\left(\rZy\right)^2\cdot\hZy} \\
   &= \rhoZ \approx \result{\rhoZgP{0}{2}}
\end{align}
Die Dichtetabelle (FoTaBe S.~187) enthält Kupfer mit \rCuO\ (und Nickel mit \rNiO). Der Zylinder besteht also $\result{\text{aus Kupfer}}$ (oder aus Nickel, das fast dieselbe Dichte hat; die Farbe unterscheidet die beiden).
